\documentclass{article}%
\usepackage[T1]{fontenc}%
\usepackage[utf8]{inputenc}%
\usepackage{lmodern}%
\usepackage{textcomp}%
\usepackage{lastpage}%
\usepackage{geometry}%
\geometry{margin=1in}%
\usepackage{graphicx}%
%
\usepackage{amsmath}%
\usepackage{amssymb}%
\usepackage{booktabs}%
\usepackage{longtable}%
\usepackage{float}%
\usepackage{caption}%
\usepackage{hyperref}%
\title{CHE 565 {-}{-} Homework 2}%
\author{Soki Sem}%
\date{\today}%
%
\begin{document}%
\normalsize%
\setlength{\parindent}{0pt}%
\setlength{\parskip}{0.6em}%
\maketitle%
\section*{Problem 1}%
\label{sec:Problem1}%
\subsection*{Setup}%
\label{subsec:Setup}%
\begin{align*}%
\text{Option A:}\quad C_{0,A} &= \$3{,}800{,}000 \\%
\quad\quad\quad\;\;\; FV_A &= \$1{,}100{,}000/\text{yr} \\%
\text{Option B:}\quad C_{0,B} &= \$5{,}000{,}000 \\%
\quad\quad\quad\;\;\; FV_B &= \$1{,}410{,}000/\text{yr}%
\end{align*}%
PV = present value, FV = future value (annual cash flow), r = yearly interest rate, n = number of years%
\par%
n= 10 years, r = 0.10%
\par

%
\subsection*{a.) Compute Net Present Value for Options A and B}%
\label{subsec:a.)ComputeNetPresentValueforOptionsAandB}%
\begin{align*}%
\text{Present Value formula:}\ \text{PV} &= \frac{F \cdot \left(r + 1\right)^{- n} \cdot \left(\left(r + 1\right)^{n} - 1\right)}{r} \\%
\text{Net Present Value formula:}\ \text{NPV} &= C_{0} + \frac{F \cdot \left(r + 1\right)^{- n} \cdot \left(\left(r + 1\right)^{n} - 1\right)}{r}%
\end{align*}%
\[%
\text{NPV}_A = \$2,959,023.82,\quad \text{NPV}_B = \$3,663,839.62%
\]%
NPV is higher for option B at 10\% yearly interest, so B is preferred under these assumptions.%
\par

%
\subsection*{b.) Annual Payment for 10 Year Lifetime, No Salvage Value, 5\% Interest Rate}%
\label{subsec:b.)AnnualPaymentfor10YearLifetime,NoSalvageValue,5InterestRate}%
P = annual payment, C0 = initial cost, r = yearly interest rate, n = number of years%
\par%
\[%
\text{Annual payment formula:}\ \text{P} = \frac{C_{0} \cdot r \cdot \left(r + 1\right)^{n}}{\left(r + 1\right)^{n} - 1}%
\]%
\[%
\text{Annual payment for option A:}\ \text{P}_A = \$492,117.38/\text{year}%
\]%
\[%
\text{Annual payment for option B:}\ \text{P}_B = \$647,522.87/\text{year}%
\]

%
\section*{Problem 2}%
\label{sec:Problem2}%
\subsection*{Setup}%
\label{subsec:Setup}%
\$0.94/lbs is the price of powdered detergent.%
\par%
0.1 \% of 4,000,000 lbs of detergent per year is carried through the exhaust%
\par%
Adding a second cyclone separator costs \$10,000, with \$300/year in maintenance costs will reduce carryover to 0.0\%, with an 8\% yearly interest rate.%
\par%
Cost of carryover detergent: \$3,760.00/year\newline%
Cost of second separator: \$10,000.00 initial + \$600.00/year maintenance = \$10,600.00 total%
\par%
Lost detergent value: \$3,760.00/year%
\par

%
\subsection*{a.) Find the additional yearly income}%
\label{subsec:a.)Findtheadditionalyearlyincome}%
Gross additional income (recovered detergent): \$3,760.00/year%
\par%
Net additional income after maintenance: \$3,460.00/year%
\par

%
\subsection*{b.) Find the payback period}%
\label{subsec:b.)Findthepaybackperiod}%
\[%
\text{The NPV equation for the payback period is:\\} -10000 + 43250.0 \cdot {1.08}^{- n} \cdot \left({1.08}^{n} - 1\right)%
\]%
Discounted payback period (NPV = 0): 3.42 years%
\par

%
\section*{Problem 3}%
\label{sec:Problem3}%
\subsection*{Setup}%
\label{subsec:Setup}%
Determine if each function is convex%
\par%
\begin{align*}%
\text{A}: f(x) &= (x_1-x_2)^2 + x_2^2 \\%
\text{B}: f(x) &= x_1^2 + x_2^2+x_3^2 \\%
\text{C}: f(x) &= e^{x_1}+e^{x_2}%
\end{align*}

%
\subsection*{Results}%
\label{subsec:Results}%
\begin{table}[H]%
\centering%
\begin{tabular}{c|c|c|c|c}%
\hline%
Function&Hessian&Eigenvalues&Convex?&Reason\\%
\hline%
$x_{2}^{2} + \left(x_{1} - x_{2}\right)^{2}$&$\left[\begin{matrix}2 & -2\\-2 & 4\end{matrix}\right]$&$\left[\begin{matrix}3 - \sqrt{5}\\\sqrt{5} + 3\end{matrix}\right]$&Yes&Eigenvalues are non{-}negative\\%
\hline%
$x_{1}^{2} + x_{2}^{2} + x_{3}^{2}$&$\left[\begin{matrix}2 & 0 & 0\\0 & 2 & 0\\0 & 0 & 2\end{matrix}\right]$&$\left[\begin{matrix}2\end{matrix}\right]$&Yes&Eigenvalues are non{-}negative\\%
\hline%
$e^{x_{1}} + e^{x_{2}}$&$\left[\begin{matrix}e^{x_{1}} & 0\\0 & e^{x_{2}}\end{matrix}\right]$&$\left[\begin{matrix}e^{x_{1}}\\e^{x_{2}}\end{matrix}\right]$&Yes&Eigenvalues are non{-}negative\\%
\hline%
\end{tabular}%
\end{table}

%
\section*{Problem 4}%
\label{sec:Problem4}%
Is the region defined by the constraints convex?%
\par%
The feasible region is defined by the following constraints: %
$x_{2} \geq 0 \wedge x_{2} \geq 1 - x_{1} \wedge x_{1} \leq 2 \wedge x_{2} \leq \frac{x_{1}}{2} + 1$%
\par%


\begin{figure}[H]%
\centering%
\includegraphics[width=0.95\textwidth]{HW2_CHE565_plot.png}%
\caption{Feasible Region}%
\end{figure}

%
The feasible region is convex because it is defined by a set of linear inequalities, which form a convex set.%
\par

%
\section*{Problem 5}%
\label{sec:Problem5}%
\subsection*{Setup}%
\label{subsec:Setup}%
Maximize the objective function:%
\par%
\[%
f(x) = - \frac{x^{6}}{6} + \frac{4 x^{5}}{5} - \frac{x^{4}}{4} - \frac{10 x^{3}}{3} + 2 x^{2} + 8 x + 1%
\]%
Given that:%
\par%
\[%
\frac{df}{dx}= \left(2 - x\right)^{3} \left(x + 1\right)^{2}%
\]%
\[%
\frac{d^2f}{dx^2}= \left(2 - x\right)^{3} \left(2 x + 2\right) - 3 \left(2 - x\right)^{2} \left(x + 1\right)^{2}%
\]

%
\subsection*{a.) Analytical Solution}%
\label{subsec:a.)AnalyticalSolution}%
Expanded form:%
\par%
\[%
\frac{df}{dx} = x^{5} - 4 x^{4} + x^{3} + 10 x^{2} - 4 x - 8%
\]%
\[%
\frac{d^2f}{dx^2} = 5 x^{4} - 16 x^{3} + 3 x^{2} + 20 x - 4%
\]%
The roots to df/dx are x* = {-}1, 2%
\par%
The corresponding function values are f(x*) = {-}2.88, 9.27%
\par%
The Hessian evaluated at the critical points are%
$\frac{d^2f}{dx^2}|_{x^*}$%
 = 0, 0, inconclusive 2nd derivative test%
\par%
The maximum value of the function is f(x*) = 9.27 at x* = 2%
\par

%
\subsection*{b.) Excel Solution}%
\label{subsec:b.)ExcelSolution}%
The Excel data was used to perform root{-}finding using the Newton{-}Raphson method:%
\par%
\[%
x_{n+1} = x_n - \frac{f'(x_n)}{f''(x_n)}%
\]%
\begin{table}[H]%
\centering%
\begin{tabular}{cccccc}%
\hline%
Iteration&x\_n&f(x\_n)&f'(x\_n)&f''(x\_n)&x\_n+1\\%
\hline%
0.0000&5.0000&{-}586.0833&{-}972.0000&{-}1296.0000&4.2500\\%
\hline%
1.0000&4.2500&{-}139.2208&{-}313.9541&{-}538.2070&3.6667\\%
\hline%
2.0000&3.6667&{-}27.0988&{-}100.8230&{-}224.6914&3.2179\\%
\hline%
3.0000&3.2179&0.5555&{-}32.1432&{-}94.4150&2.8775\\%
\hline%
4.0000&2.8775&7.2322&{-}10.1590&{-}39.9713&2.6233\\%
\hline%
5.0000&2.6233&8.8043&{-}3.1799&{-}17.0590&2.4369\\%
\hline%
6.0000&2.4369&9.1644&{-}0.9854&{-}7.3392&2.3027\\%
\hline%
7.0000&2.3027&9.2446&{-}0.3025&{-}3.1810&2.2076\\%
\hline%
8.0000&2.2076&9.2620&{-}0.0920&{-}1.3876&2.1413\\%
\hline%
9.0000&2.1413&9.2657&{-}0.0278&{-}0.6084&2.0955\\%
\hline%
10.0000&2.0955&9.2665&{-}0.0084&{-}0.2678&2.0643\\%
\hline%
11.0000&2.0643&9.2666&{-}0.0025&{-}0.1182&2.0432\\%
\hline%
12.0000&2.0432&9.2667&{-}0.0007&{-}0.0523&2.0289\\%
\hline%
13.0000&2.0289&9.2667&{-}0.0002&{-}0.0232&2.0193\\%
\hline%
14.0000&2.0193&9.2667&{-}0.0001&{-}0.0103&2.0129\\%
\hline%
15.0000&2.0129&9.2667&{-}0.0000&{-}0.0046&2.0086\\%
\hline%
16.0000&2.0086&9.2667&{-}0.0000&{-}0.0020&2.0058\\%
\hline%
17.0000&2.0058&9.2667&{-}0.0000&{-}0.0009&2.0038\\%
\hline%
18.0000&2.0038&9.2667&{-}0.0000&{-}0.0004&2.0026\\%
\hline%
19.0000&2.0026&9.2667&{-}0.0000&{-}0.0002&2.0017\\%
\hline%
20.0000&2.0017&9.2667&{-}0.0000&{-}0.0001&2.0011\\%
\hline%
\end{tabular}%
\caption{Optimization Results from Excel starting at x0=5}%
\end{table}%
\begin{table}[H]%
\centering%
\begin{tabular}{cccccc}%
\hline%
Iteration&x\_n&f(x\_n)&f'(x\_n)&f''(x\_n)&x\_n+1\\%
\hline%
0.0000&{-}5.0000&{-}4832.7500&5488.0000&{-}5096.0000&{-}3.9231\\%
\hline%
1.0000&{-}3.9231&{-}1208.5512&1775.5079&{-}2114.1044&{-}3.0832\\%
\hline%
2.0000&{-}3.0832&{-}295.6352&570.0315&{-}883.6738&{-}2.4382\\%
\hline%
3.0000&{-}2.4382&{-}71.0801&180.8132&{-}373.6709&{-}1.9543\\%
\hline%
4.0000&{-}1.9543&{-}17.8527&56.3066&{-}160.7261&{-}1.6040\\%
\hline%
5.0000&{-}1.6040&{-}5.9169&17.0747&{-}70.7559&{-}1.3626\\%
\hline%
6.0000&{-}1.3626&{-}3.4409&5.0003&{-}32.0380&{-}1.2066\\%
\hline%
7.0000&{-}1.2066&{-}2.9756&1.4068&{-}14.9373&{-}1.1124\\%
\hline%
8.0000&{-}1.1124&{-}2.8972&0.3808&{-}7.1437&{-}1.0591\\%
\hline%
9.0000&{-}1.0591&{-}2.8853&0.0999&{-}3.4805&{-}1.0304\\%
\hline%
10.0000&{-}1.0304&{-}2.8836&0.0257&{-}1.7158&{-}1.0154\\%
\hline%
11.0000&{-}1.0154&{-}2.8834&0.0065&{-}0.8515&{-}1.0078\\%
\hline%
12.0000&{-}1.0078&{-}2.8833&0.0016&{-}0.4242&{-}1.0039\\%
\hline%
13.0000&{-}1.0039&{-}2.8833&0.0004&{-}0.2117&{-}1.0020\\%
\hline%
14.0000&{-}1.0020&{-}2.8833&0.0001&{-}0.1057&{-}1.0010\\%
\hline%
15.0000&{-}1.0010&{-}2.8833&0.0000&{-}0.0528&{-}1.0005\\%
\hline%
16.0000&{-}1.0005&{-}2.8833&0.0000&{-}0.0264&{-}1.0002\\%
\hline%
17.0000&{-}1.0002&{-}2.8833&0.0000&{-}0.0132&{-}1.0001\\%
\hline%
18.0000&{-}1.0001&{-}2.8833&0.0000&{-}0.0066&{-}1.0001\\%
\hline%
19.0000&{-}1.0001&{-}2.8833&0.0000&{-}0.0033&{-}1.0000\\%
\hline%
20.0000&{-}1.0000&{-}2.8833&0.0000&{-}0.0017&{-}1.0000\\%
\hline%
\end{tabular}%
\caption{Optimization Results from Excel starting at x0={-}5}%
\end{table}%
\begin{table}[H]%
\centering%
\begin{tabular}{ccc}%
\hline%
x0&x*&f(x*)\\%
\hline%
5.0000&2.0011&9.2667\\%
\hline%
{-}5.0000&{-}1.0000&{-}2.8833\\%
\hline%
\end{tabular}%
\caption{Summary of Optimization Results from Excel}%
\end{table}

%
\subsection*{c.) Numerical solution using Nelder{-}Mead Simplex algorithm which is the same as fminsearch in matlab}%
\label{subsec:c.)NumericalsolutionusingNelder{-}MeadSimplexalgorithmwhichisthesameasfminsearchinmatlab}%
Using Nelder{-}Mead Simplex algorithm with initial guesses x0 = 5: %
\par%
One of the roots of df/dx are found to be x* = 1.9999%
\par%
The corresponding optimum function values are f(x*) = 9.2667%
\par%
Actual maximum of the function:%
\par%
\begin{align*}%
x^* &=  1.9999 \\%
f(x^*) &= 9.2667 \\%
\text{iterations} &= 20 \\%
\text{function calls} &= 41%
\end{align*}

%
\subsection*{d.) Numerical solution using Newton Conjugate Gradient Method which is the same as fmincon in matlab except required to give jacobian and hessian information}%
\label{subsec:d.)NumericalsolutionusingNewtonConjugateGradientMethodwhichisthesameasfminconinmatlabexceptrequiredtogivejacobianandhessianinformation}%
Using the Newton Conjugate Gradient method with initial guesses x0 = 5: %
\par%
One of the roots of df/dx are found to be x* = 2.0038%
\par%
The corresponding optimum function values are f(x*) = 9.2667%
\par%
\begin{align*}%
x^* &= 2.0038 \\%
f(x^*) &= 9.2667 \\%
\text{iterations} &= 19 \\%
\text{function calls} &= 19%
\end{align*}%
Using the Newton Conjugate Gradient method with initial guesses x0 = {-}5: %
\par%
One of the roots of df/dx are found to be x* = {-}1.0001%
\par%
The corresponding optimum function values are f(x*) = {-}2.8833%
\par%
\begin{align*}%
x^* &= -1.0001 \\%
f(x^*) &= -2.8833 \\%
\text{iterations} &= 20 \\%
\text{function calls} &= 20%
\end{align*}

%
\subsection*{e.) Compare results from all methods}%
\label{subsec:e.)Compareresultsfromallmethods}%
\begin{table}[H]%
\centering%
\begin{tabular}{cccccc}%
\hline%
Method&Initial Guess&x*&f(x*)&Iterations&Function Calls\\%
\hline%
Nelder{-}Mead&5&1.9999&9.2667&20&41\\%
\hline%
Excel&{-}5&2.0011&9.2667&$21$&$21$\\%
\hline%
Excel&5&{-}1.0000&{-}2.8833&$21$&$21$\\%
\hline%
Newton{-}CG&5&2.0038&9.2667&19&19\\%
\hline%
Newton{-}CG&{-}5&{-}1.0001&{-}2.8833&20&20\\%
\hline%
\end{tabular}%
\end{table}%
All methods found the same maximum value of the function at x* = 2 and f(x*) = 9.267, which is consistent with the analytical solution. The Nelder{-}Mead method required 20 iterations and 41 function calls, the Newton{-}CG method converged in 20 iterations, the Excel solution also converged to the same optimum value, demonstrating consistency across methods. However, Newton{-}CG and the excel solution are significantly more influenced by the initial guess, and can converge to a local extremum. In the case of this polynomial function, there are two critical points based on the roots of the derivative, one is an optimum at x* = 2 and the other is a saddle point at x* = {-}1, and the initial guess determines which extremum the algorithm converges to. The Nelder{-}Mead method is less sensitive to the initial guess and was able to find the global maximum regardless of the starting point.%
\par

%
\end{document}